\subsection{Limitations of the Evidence-Synthesis Process}

The conclusions of this study are constrained first by the design of the
evidence synthesis itself. The literature was identified through predefined
translational research questions spanning memory neuroscience, sleep
physiology, TMR, wearable sensing, real-time state estimation, and closed-loop
intervention. This structure enabled evidence from otherwise separate fields to
be evaluated within a common development problem, but it was not designed to
provide exhaustive retrieval of every publication within each domain.

The resulting evidence base may therefore omit relevant work because of
differences in terminology, indexing, disciplinary boundaries, publication
access, or the progressive refinement of the research questions. This affects
the comprehensiveness and reproducibility of the synthesis more directly than
the validity of any single cited experiment. The pathway selected in Section~6
should consequently remain revisable if additional evidence materially changes
the relative support for competing implementation strategies.

A second limitation is reviewer-dependent judgment. Source identification,
critical appraisal, and translational interpretation were conducted by a single
reviewer without duplicate screening, independent extraction, formal
risk-of-bias scoring, or standardized certainty grading. The predefined
research-question structure, explicit source-function criteria, preference for
primary and independently validated evidence, and claim-level traceability
described in Section~2 provided safeguards against unrestricted source
selection, but they do not eliminate subjective judgment.

No new quantitative synthesis was performed. Effect estimates were taken from
published meta-analyses or individual studies rather than recomputed across the
assembled literature. The present study therefore cannot provide a new pooled
estimate of TMR efficacy, formally rank the certainty of competing sensing
pathways, or assign quantitative confidence to the pathway comparison. Its
outputs should be interpreted as evidence-supported translational judgments
rather than as systematic-review certainty estimates or independently pooled
effect sizes.


\subsection{Limitations of the Available Scientific Evidence}

The TMR literature itself places important limits on translation. Experimental
protocols vary in memory domain, initial learning strength, cue modality,
stimulus intensity, sleep stage, cue timing, physiological monitoring,
participant characteristics, retention interval, and behavioral outcome.
Accordingly, TMR should not be treated as a single standardized intervention.
The meta-analysis by Hu et al.\ identified a positive average effect
(\(g=0.29\)) together with substantial between-study heterogeneity
(\(I^2=71\%\)) \citep{hu2020promoting}. The same analysis identified funnel-plot
asymmetry consistent with possible publication bias; under one trim-and-fill
analysis, the estimated effect was reduced while remaining positive
\citep{hu2020promoting}. The pooled value should therefore be interpreted as an
estimate across heterogeneous experimental conditions rather than as an
expected performance target for an engineered system.

External validity is also limited. Much of the evidence relevant to TMR derives
from healthy participants completing tightly controlled memory tasks over one
or a small number of nights. Such experiments are well suited to identifying
causal effects, but they provide much less information about repeated
unsupervised use, night-to-night reliability, long retention intervals,
habituation to cues, adherence, age-related variation, atypical sleep
architecture, or clinical populations. Existing home studies extend ecological
validity but remain limited in number, scale, duration, and task diversity
\citep{goldi2019home,whitmore2022improving,salfi2025promoting}.

The long-term behavioral and physiological consequences of repeated TMR also
remain insufficiently characterized. Short experimental protocols can determine
whether appropriately delivered cues alter subsequent memory without
establishing whether repeated intervention across weeks or months remains
effective, produces habituation, changes normal sleep architecture, or affects
memory processes beyond the intended material. The present evidence does not
demonstrate such adverse effects, but neither does it justify treating their
absence in short studies as evidence of long-term safety. Claims concerning
durability, cumulative benefit, and repeated unsupervised exposure must
therefore await longitudinal evaluation.

Mechanistic interpretation remains incomplete as well. Cue-evoked spindle
activity, content-related EEG patterns, and temporal sensitivity to endogenous
oscillatory dynamics are consistent with contemporary models of memory
reactivation, but no single neurophysiological marker has been established as a
validated surrogate for behavioral TMR success. Consequently, a system should
not infer successful memory enhancement solely from detection of a spindle,
slow wave, classifier output, or other physiological response.

The wearable literature introduces a further evidence boundary. Studies of
wearable EEG, PPG, and consumer sleep-tracking systems differ substantially in
population, hardware, signal preprocessing, class definitions, reference
scoring, algorithm version, and validation setting
\citep{arnal2020dreem,korkalainen2020ppg,chinoy2021consumer,
ravindran2025dreem}. Performance reported for one device or algorithm cannot
therefore be generalized to an entire sensing modality. Furthermore, much of
the wearable evidence concerns retrospective agreement with PSG rather than
causal operation within an intervention loop. Real-time TMR requires additional
properties---including usable continuous signals, uncertainty handling,
synchronization, bounded sensing-to-cue delay, abstention during unreliable
intervals, and post-cue monitoring---that are rarely validated together in the
same system.


\subsection{Limitations of Translational Inference}

The central translational limitation is that evidence for individual components
does not constitute validation of their integration. The preceding sections
provide support for several separate propositions: TMR can produce a modest
memory-specific behavioral effect; sleep physiology can be observed through
EEG and estimated through peripheral signals; reduced wearable systems can
operate outside full PSG; and both EEG- and peripheral-guided systems have been
used to automate cue delivery at home
\citep{whitmore2022improving,salfi2025promoting}. None of these findings
individually demonstrates that the first-generation configuration selected in
Section~6 will operate reliably as an integrated system.

Integration creates new failure modes at the interfaces among sensing,
preprocessing, state estimation, control logic, communication, cue execution,
and post-stimulation monitoring. A physiological classifier may perform
adequately while the end-to-end intervention fails because of signal loss,
timing error, synchronization failure, inappropriate thresholds, or an
unrecognized post-cue transition. Conversely, a technically reliable
closed-loop implementation may operate exactly as intended without producing a
measurable behavioral benefit. Technical validity and biological efficacy must
therefore be demonstrated separately before they can be interpreted jointly.

The EEG-first decision is also not the result of direct comparative
validation. No evidence reviewed in this study establishes that an EEG-guided
controller is behaviorally superior to an otherwise equivalent PPG- or
multimodal-peripheral controller under matched participants, learning tasks,
cue policies, and physiological conditions. Similarly, the evidence does not
establish that stage-aware stimulation is superior to phase-aware control, or
that integration with existing hardware is superior to proprietary hardware in
all relevant dimensions.

The selected reference pathway is consequently a sequencing decision based on
evidentiary continuity and the number of additional assumptions introduced into
the initial validation problem. EEG was selected because it preserves closer
access to the electrophysiological domain used to characterize sleep and has
already supported automated home TMR, not because permanent superiority over
peripheral sensing has been demonstrated. Stage-aware control was selected
because its behavioral evidence base is broader and its implementation requires
fewer timing assumptions than phase-locked control, not because it has been
shown to be physiologically optimal.

The translation from scientific evidence to system requirements is itself
qualitative. Requirements such as physiological eligibility, conservative
abstention, temporal validity, and post-cue monitoring follow from the reviewed
evidence, but the literature does not provide universal numerical thresholds
for implementing them. The present analysis therefore identifies
\textit{what must be validated} more strongly than it determines
\textit{the numerical values at which validation will succeed}.


\subsection{Boundaries of the Present Conclusions}

Within these limitations, the evidence supports controlled development and
prospective evaluation of an EEG-guided, stage-aware, real-time closed-loop TMR
reference pathway. It does not establish that the resulting system will produce
reliable memory benefits across participants, tasks, recording nights, devices,
or broader populations. It also does not establish that an integrated system
can yet maintain sufficient signal quality, state-estimation reliability,
timing accuracy, disturbance control, or behavioral efficacy during extended
unsupervised operation.

The present study does not provide validated numerical thresholds for sleep-state
confidence, signal quality, cue intensity, cue spacing, state stability,
acceptable latency, or post-stimulation arousal. These variables should remain
configurable during experimental development and should be determined through
platform-specific physiological and behavioral validation rather than fixed
from indirect literature estimates.

Nor does the paper establish a finalized system architecture, production
hardware configuration, deployable software product, clinical intervention, or
consumer-ready memory-enhancement technology. The selected configuration is a
\textit{first-generation reference pathway}: a testable arrangement chosen to
provide an interpretable starting point for engineering and experimental
validation.

The conclusions can therefore be stated at four different levels. The
\textbf{biological evidence} supports TMR as a condition-dependent means of
influencing sleep-dependent memory processing. The \textbf{experimental
evidence} supports a modest average behavioral effect under appropriate
conditions. The \textbf{translational analysis} supports EEG-guided,
stage-aware closed-loop control as the most defensible initial reference
pathway under the evidence reviewed. The \textbf{product-level evidence},
however, remains insufficient because the proposed integrated system has not
yet been implemented and prospectively validated.

These boundaries are not peripheral qualifications to the central conclusion;
they define its valid scope. The present work narrows the development space,
identifies the physiological and measurement functions that should be preserved,
and establishes a reference configuration against which subsequent engineering
choices can be tested. The remaining uncertainties therefore become explicit
research questions rather than hidden assumptions, providing the basis for the
validation priorities described in Section~8.